\documentclass[11pt]{article}
\usepackage[a4paper,margin=25mm]{geometry}
\usepackage{amsmath,amssymb,amsthm,booktabs,hyperref,xurl}
\newtheorem{theorem}{Theorem}
\newtheorem{lemma}[theorem]{Lemma}
\newtheorem{proposition}[theorem]{Proposition}
\newcommand{\E}{\mathrm E}
\newcommand{\D}{\mathrm D}
\newcommand{\A}{\mathrm A}
\title{Triangle constraints at two compact arithmetic rank--21 vertex types}
\author{Research manuscript with reproducible exact certificates}
\date{18 October 2026}
\begin{document}
\maketitle
\begin{abstract}
We prove local incidence restrictions at compact arithmetic hyperbolic
Coxeter vertices with ground field $\mathbb Q(\sqrt5)$ and diagram
$\A_1+\D_5+\E_7+\E_8$ or $\D_5+2\E_8$.
The first type has at most seven triangular two-faces marked by
nonadjacent roots within one nontrivial component. At either type, at
most one of the three possible nonadjacent $\D_5$ triangles occurs.
We describe forced simplex and product faces in the exceptional
components. Exact finite-corner and arithmetic computations are supplied.
These results do not improve a dimension ceiling or establish a new
global charging inequality.
\end{abstract}

\section{Statements and hypotheses}
Roots have square two. Number $\D_5$ by the edges $01,12,23,24$.
Number $\E_n$ by the path $0-1-\cdots-(n-2)$ and the extra edge
$2(n-1)$. All these bonds have label three.
A pair of incident facet roots specifies a two-face at a simple vertex.
Call the pair nonadjacent when it is not an edge of its old component
diagram. We consider only pairs in a single component.

\begin{theorem}\label{thm:seven}
Let $P\subset\mathbb H^{21}$ be a compact arithmetic Coxeter polytope
with ground field $\mathbb Q(\sqrt5)$, and let a vertex have diagram
$\A_1+\D_5+\E_7+\E_8$. Among its $42$ two-faces marked by nonadjacent
pairs within $\D_5,\E_7,\E_8$, at most seven are triangles.
Equality in the necessary local domain requires the $\E_7$ third wall
with label-five contact at root~5 and the two $\E_8$ third walls with
label-five contacts at roots~0 and~6. It has no nonadjacent $\D_5$
triangle. Realization of this necessary equality configuration is not asserted.
\end{theorem}

\begin{theorem}\label{thm:d5}
At either $\A_1+\D_5+\E_7+\E_8$ or $\D_5+2\E_8$, under the same
compact arithmetic and ground-field hypotheses, at most one of the
$\D_5$ faces marked $(0,2),(1,3),(1,4)$ is triangular. Every other
nonadjacent $\D_5$ face is nontriangular.
\end{theorem}

The repository's separate, completed internal compact-dimension--21
field reduction leaves precisely this ground field. That reduction is
an input when applying these conditional statements to an arbitrary
compact arithmetic chamber in dimension~21. We do not extend the
three-field theorem valid from dimension~22 to dimension~21.

\section{Actual corners and a complete finite domain}
Let $T$ be the positive span of the $21$ old vertex roots, and let $e$
be its future unit timelike orthogonal vector. If the face marked $a,b$
is a triangle with third wall $c$, both root sets replacing $a$ or $b$
by $c$ are positive-definite actual corners.

\begin{lemma}\label{lem:isolation}
A whole untouched $\E_8$ is orthogonal to $c$. A whole untouched
$\E_7$ is also orthogonal to $c$.
\end{lemma}
\begin{proof}
Finite $\E_8$ has no connected finite extension. If $c$ meets a whole
$\E_7$, its finite extension must be $\E_8$. Maximality then forces
$c$ orthogonal to every other root in either new corner. Their union
contains all old roots. The full old forest together with $c$ would
therefore be positive definite of rank $22$, impossible in signature
$(21,1)$.
\end{proof}

For $\D_5$ marks the remaining complete old domain is $\D_5+\A_1$,
or just $\D_5$ at the second vertex type. For $\E_7$ or $\E_8$ marks
in Theorem~\ref{thm:seven}, we retain respectively the full forests
$\E_7+\D_5+\A_1$ and $\E_8+\D_5+\A_1$. No $\D_5$ or $\A_1$
contact is assumed zero in these latter enumerations.

\begin{lemma}\label{lem:labels}
Every contact in these domains has label at most five.
\end{lemma}
\begin{proof}
A label at least six in a finite Coxeter diagram belongs to an isolated
$I_2(m)$. At an old root in a nontrivial connected component this
contradicts the retained old neighbors in the two corners: the marks
are distinct and nonadjacent. If instead the whole isolated $\A_1$
has such a contact, the isolated-$I_2(m)$ requirement forces $c$
orthogonal to every other old root in both corners, hence their union.
Its height then is
\[
 H=\frac{(2\cos(\pi/m))^2}{2}-2<0,
\]
contradicting $c=p+\sqrt H e$ with positive height.
\end{proof}

The first finite corner can consequently be enumerated through finite
trees. The new root has degree at most three, meets each old connected
component at most once, and with degree three has only label-three
bonds. Restore the omitted marked root with labels two through five
and test the second corner. These conditions contain every actual
triangle, not merely a selected finite witness.

\begin{center}
\begin{tabular}{lrrr}
\toprule
Marks & First-corner trials & Second-corner trials & Finite eligible\\
\midrule
$\E_7$ in $\E_7+\D_5+\A_1$ &7497&2296&379\\
$\E_8$ in $\E_8+\D_5+\A_1$ &12492&3368&470\\
\bottomrule
\end{tabular}
\end{center}
For $\D_5+\A_1$, a separate direct enumeration checks all
$6\cdot4^6=24576$ contact vectors and has $139$ finite-eligible occurrences.

\section{Arithmetic and the surviving normals}
Put $s=\sqrt5$ and $\phi=(1+s)/2$. The computation uses exact rational
arithmetic in $\mathbb Q(\sqrt2,\sqrt5)$, the field containing every
contact allowed by Lemma~\ref{lem:labels}. For each old Gram $G$,
\[
 H=bG^{-1}b^{\mathsf T}-2.
\]
It rejects $H\le0$. Simple cycles use $c$ and the unique old path
between two contacts; the old path contributes rational factors.
Their products must belong to $\mathbb Q(s)$. Once that holds,
$H\in\mathbb Q(s)$, and conjugate positive semidefiniteness requires
$\overline H\le0$. The old rational positive Gram is unchanged, so
this follows directly from its Schur complement. Connected affine
principal subsets containing $c$ are excluded by compactness: their
positive null kernels give a future ideal point in the chamber closure.
Finite subsets, and connected subsets with label five, cannot be affine.

The $\D_5+\A_1$ domain leaves six roots:
\begin{center}\small
\begin{tabular}{llll}
\toprule
Name & $\D_5$ contacts & $\A_1$ contact & $H$\\
\midrule
$C_0$&$\phi$ at0&0&$(s-1)/2$\\
$L_3,L_4$&$\phi$ at3 or4&0&$(-1+5s)/8$\\
$U_3,U_4$&1 at1, $\phi$ at3 or4&0&$(23+13s)/8$\\
$C_*$&1 at0&$\phi$&$(s-1)/4$\\
\bottomrule
\end{tabular}
\end{center}
The first and last correspond to marks $(0,2)$; the two fork families
correspond to $(1,3)$ and $(1,4)$. In $\D_5+2\E_8$ there is no
$\A_1$, and only the first five roots occur.

With the whole $\D_5+\A_1$ included, the exceptional domains leave:
\begin{center}\small
\begin{tabular}{llll}
\toprule
Name & Exceptional contact & $\A_1$ contact & $H$\\
\midrule
$E_0$&$\E_7$: $\phi$ at0&0&$1+s$\\
$E_5$&$\E_7$: $\phi$ at5&0&$(1+3s)/4$\\
$E_*$&$\E_7$: 1 at5&$\phi$&$(1+s)/4$\\
$F_0$&$\E_8$: $\phi$ at0&0&$4+2s$\\
$F_6$&$\E_8$: $\phi$ at6&0&$1+s$\\
\bottomrule
\end{tabular}
\end{center}
All have zero $\D_5$ contact; this is an output of the complete domains.
Their nonadjacent triangle marks are respectively
\[
\begin{gathered}
E_0:(0,2),\quad E_5:(2,4),(2,5),(3,5),\quad E_*:(3,5),\\
F_0:(0,2),\quad F_6:(3,5),(3,6),(4,6).
\end{gathered}
\]

\section{Coupling at the common vertex}
For two candidate walls,
\[
 c_i\cdot c_j=p_i\cdot p_j-\sqrt{H_iH_j}.
\]
The old projections of roots from different components are orthogonal,
except for their explicitly shared $\A_1$ contacts. The ring of integers
is $\mathbb Z[\phi]$: $a+bs$ is integral exactly when $2b$ and $a-b$
are integers. With zero projection product, $H_iH_j$ is the length-two
cyclic product and must be integral.

\begin{proof}[Proof of Theorem~\ref{thm:d5}]
The exact $\D_5+\A_1$ inverse gives positive projection product $p$
and $p^2-H_iH_j>0$ for all fifteen pairs of the six candidate walls.
Thus each mutual product would be positive, contrary to acuteness.
For the new $C_*$ against $C_0,L_j,U_j$, these positive differences are
\[
 \frac{3(1+s)}4,\qquad\frac{-7+5s}{16},\qquad\frac{9+5s}{16}.
\]
The ten pairs among the old five are independently checked as well.
Distinct marked faces use disjoint candidate sets, proving the result.
\end{proof}

\begin{proof}[Proof of Theorem~\ref{thm:seven}]
The fifty-five candidate pairs split into eighteen forbidden positive
mutual products, twenty-three nonintegral length-two cycles, two
orthogonal necessary pairs, and twelve remaining necessary native or
dotted possibilities. The shared $\A_1$ is essential: for $C_*,E_*$
their actual product is $(1+s)/4>0$.

Within $\E_7$, $E_5$ is incompatible with $E_0,E_*$, whereas $E_0,E_*$
are orthogonal. Thus at most three designated triangles occur there.
The two $\E_8$ roots are orthogonal and supply at most four. If any
$\D_5$ triangle occurs, every $\E_7$ root except possibly $E_0$ is
forbidden, giving at most $1+1+4=6$. Otherwise the total is at most
$3+4=7$. The exact checker also enumerates all $2^{11}=2048$ necessary
root subsets and gives the unique maximizing subset
$\{E_5,F_0,F_6\}$.
\end{proof}

\section{Forced exceptional faces and the remaining charge problem}
An actual $\E_8$ wall $F_0$ forces all three triangles on marks
$\{0,1,2\}$, and $F_6$ forces all six on $\{3,4,5,6\}$. The checked
positive corners are actual vertices by projection of a strict interior
timelike point off a positive Gram span. The inverse is nonnegative,
so all other inequalities remain strict. The resulting simplex-edge
graphs saturate the degree of each simple face and hence exhaust its
connected edge graph. These are actual simplex three- and four-faces.

The inverse entries $J_{00}=4,J_{66}=2,J_{06}=2$ give
\[
 p_0\cdot p_6=2\phi^2=3+s,
 \qquad H_0H_6=(3+s)^2,
\]
so the walls are orthogonal. If both occur, all twenty corners obtained
by omitting one member of $\{0,1,2,F_0\}$ and one of
$\{3,4,5,6,F_6\}$ are positive. Keeping root7 and the other old
components as supports gives a seven-face with the full incidence
structure $\Delta^3\times\Delta^4$, again by degree saturation and
connectedness. Its intrinsic dihedral cosine at old roots1,2 is
$1/\sqrt3$, so it is not intrinsically Coxeter.

For $\D_5+2\E_8$, the earlier complete nonadjacent domain has an upper
bound nine on its forty-eight designated triangles. The present
$\D_5$ strengthening does not reduce it to eight: one long $\D_5$
triangle and four from each $\E_8$ remain in the necessary domain.

Finally, a nontriangle can be a quadrilateral. The current global rule
requires total payment at least two on a triangle and at least one on
a quadrilateral. The results here provide no new quadrilateral payment,
vertex-allowance saving, dimension--20 theorem, or example. Their next
use requires actual quadrilateral or edge compatibility and a verified
change to the global charge system.

\section{Reproducibility and verification scope}
The accompanying proof map is
\path{docs/compact21_supported_triangle_constraints_map.md}.
The two complete new domain producers support \texttt{--verify};
\path{check_compact21_e8_triangle_bundles_20261018.py} is a
standalone exact arithmetic check of the earlier $\D_5$ and $\E_8$
identities and all twenty product-face corners. Sources and field
reduction inputs are hash-bound. Generation and separate full replays
pass with measured peak resident memory below $32$ MiB. The jobs use
one global lock, a $59$-second alarm and a periodic one-GiB memory guard.
No optimizer or Lean process is launched.

The universal actual-corner, compactness, projection and edge-graph
arguments are written proofs. Published finite/affine classifications,
compact simplicity and Vinberg's arithmetic criterion are explicit
theoretical inputs. The earlier $\D_5+2\E_8$ domain and complete
dimension--21 field reduction are separate internally checked inputs.
No end-to-end Lean or external human verification is claimed.
\end{document}
